\subsection{Overview}
Reducing nitrate--nitrogen (NO$_3$-N) losses to subsurface (tile) drainage and the underlying groundwater remains one of the least settled objectives in drainage science: although a wide portfolio of management levers has been proposed, the field is still dominated by open questions about which practices work, under what conditions, at what cost, and with what unintended tradeoffs. Across sites from Iowa, Illinois, Indiana, Minnesota, and the Dakotas in the U.S.\ Midwest to Canada, and across Europe (Finland, Denmark, Sweden, Norway, France, Germany, Switzerland, Italy, Ireland, Latvia) and beyond, authors repeatedly frame their contributions less as definitive answers than as calls for further field experiments, better process understanding, and improved decision-support tools. \textcite{ROSSI1991} state the underlying urgency plainly: management directed at substantial nitrate reductions is needed, together with careful evaluation of the effects of fertilizer timing, rate, form, and placement on NO$_3^-$ loss via subsurface drainage.

Several broad lines of evidence and research need recur throughout this theme. First, \emph{fertilizer N management}---rate, timing, form, placement, and crediting of residual and mineralizable N---is widely seen as the primary but incompletely optimized lever. Second, \emph{cropping-system change}, especially cover crops, perennials, and alternative rotations, is promoted but still under-quantified across soils and climates. Third, \emph{drainage water management} (controlled drainage, shallow drainage, water-table control) and \emph{edge-of-field treatment} (bioreactors, saturated riparian buffers, wetlands) are promising but require optimization and long-term verification. Fourth, \emph{process-based modeling} is expected to carry much of the predictive burden as concerns about nitrate export grow, yet substantial algorithmic gaps remain---particularly for winter, soil C--N cycling, and groundwater interactions. Fifth, \emph{monitoring, tracers, and scaling} methods must improve to apportion sources and extend site findings to landscapes. Finally, cutting across all of these are unresolved \emph{tradeoffs}---between nitrate reduction and yield, greenhouse-gas emissions, phosphorus loss, and economic risk---that the field has not yet adequately characterized. \textcite{Du2017} anticipate that models will be relied upon extensively to estimate management effects as concern over nitrate export increases, and can enhance understanding of how N management affects water quality.

\subsection{Fertilizer nitrogen management: rate, timing, form, and crediting}
The most consistently identified priority is refining fertilizer N management. \textcite{Zhao2000} argue that efforts should concentrate on identifying N management strategies that increase N-use efficiency in corn production, pointing specifically to optimal rate, timing, and form as targets for reducing nitrate levels in the Minnesota and upper Mississippi rivers. \textcite{Randall2001} similarly call to fine-tune fertilizer N management by applying the correct rate at the optimum time and crediting previous legume crops and animal-manure applications, and to improve animal-manure management by knowing nutrient content and application rate, spreading uniformly, and incorporating in a timely manner. \textcite{Emmett2022} note that research determining optimal rates and timing of N fertilization can help reduce N$_2$O emissions and minimize tradeoffs between yield and environmental performance. At the watershed scale, \textcite{Dattamudi2020} conclude that N application schemes need reevaluation for better N management and environmental sustainability, suggesting lower fertilizer rates for corn, precision agriculture, use of soil N with high-yielding varieties, split fall/summer application, and proper timing. \textcite{Elmi2002} add that reductions may be possible by carefully managing fertilization rates and timing to match crop uptake.

\subsubsection{Crediting residual, mineralizable, and legacy nitrogen}
A recurring, more specific gap is the crediting of N supplied by soil organic matter, prior legumes, manure, and land-use legacies when setting application rates. \textcite{Bakhsh2007} state a need to assess proper N credit due to mineralization from the soil organic-matter pool and soybean residue when deciding N application rates from organic or inorganic sources, and \textcite{Bakhsh2005} likewise call for assessing proper N credit to soybean and the storage of mineralizable N in the root zone before applications, to reduce NO$_3$-N leaching potential with tile flows. \textcite{Randall2001} extend this to a methodological need: improved soil-N testing methods to estimate mineralizable N and carryover N, especially after dry years, legumes, or manure applications. In systems transitioning from perennial cover, \textcite{Huggins2001} propose adjusting N recommendations to include second-year credits for N mineralization following perennial-grass conversion to row crops, note that N strategies for corn must be developed to efficiently use the release of mineralized N from CRP land, and observe that conservation tillage during CRP/alfalfa conversion would likely slow mineralization and alter residual soil N buildup and NO$_3$-N losses. \textcite{Weed1996} recommend minimizing NH$_3$ fertilizer applied to continuous- and rotation-corn plots by using late-spring soil-nitrate tests. On land plowed out of alfalfa, \textcite{ROBBINS1980} counsel care in determining N rates and timing and note that improved timing and amount of applied water could decrease required N rates while making more N available to the crop.

For irrigated systems on vulnerable soils, \textcite{Derby2009} enumerate management options---reducing N application during the first years after irrigation initiation, setting reasonable yield goals, giving proper credit for fall residual soil NO$_3$-N, and adjusting applications for adverse weather or poor soil properties---and emphasize that field-scale studies spanning soil, vadose-zone, and groundwater levels across multiple management periods are necessary to understand water quality under irrigated agriculture. \textcite{Hurtado2011} report that lower N rates in cotton could reduce losses during the cotton season and the beginning of the following one. \textcite{Ale2012b} caution that model-predicted excess soil-N reserves at the studied N rates must be verified, while noting the potential to further reduce fertilizer from the lowest rates given soil-N budget surpluses, which may further cut NO$_3$-N losses toward discharge-reduction goals. The persistence and role of \emph{legacy} N is itself flagged as underexplored: \textcite{Williams2022} argue that a clearer definition of legacy nutrients is urgently needed to better apportion losses and develop mitigation strategies, and \textcite{Hu2024} call for research into the coupled biogeochemical and hydrological processes that regulate the formation of legacy N stores and their impact on NO$_3^-$ export in tile-drained agroecosystems.

\subsubsection{Enhanced-efficiency fertilizers, stabilizers, and application method}
Beyond rate and timing, several authors point to product- and method-level questions. \textcite{Kaur2024c} recommend that future research address enhanced-efficiency fertilizer technology---spring-applied anhydrous ammonia with a nitrification inhibitor---in conjunction with enhanced drainage water management in corn--soybean rotations. \textcite{Vincent2025} note that management of NH$_4^+$-N specifically can be achieved through N stabilizers applied alongside fertilizers, and that the effects of cover crops on field- and watershed-scale NH$_4^+$-N dynamics deserve further exploration. Regarding manure delivery, \textcite{Gamble2018} report that drainage NO$_3$-N concentrations from fertigated silage corn need further study to determine whether they can be reduced to acceptable levels without impacting silage yield, and stress that tradeoffs of fertigation---crop injury, disease, pathogen transport to waterways, and alternative N-loss pathways relative to slurry injection---require further exploration. \textcite{Bechmann2014} recommend developing a strategy to even out manure application within catchments of high livestock density, and that mitigation in cereal areas should focus on reducing soil mineral N in spring and autumn.

\subsection{Timing of loss: the early-spring and seasonal window}
A distinctive and repeatedly emphasized frontier is the temporal concentration of losses, particularly in early spring, and the need for management targeted to that window. \textcite{Bakhsh2007} state that strategies need to be developed to reduce early-spring NO$_3$-N leaching losses, and \textcite{Bakhsh2010} extend this to reducing off-site nitrate transport during the March-through-June period, adding that the temporal distribution of rainfall during the growing season needs further attention when evaluating and managing nitrate leaching to drainage water. \textcite{Nance2007} argue that assessment of soil NO$_3$-N status, tile flow, and NO$_3$-N loss during April, May, and June warrants further study. \textcite{Gentry2024} suggest that a reliable, winter-hardy cover crop ahead of corn may help reduce tile nitrate in corn--soybean rotations, and that moving fall N application to spring while using overwintering grass cover crops may be the best way to decrease N loss in intensively tile-drained watersheds. For drainage-outlet management, \textcite{Skaggs1995} recommend selectively managing drainage outlets to reduce peak NO$_3$-N losses, especially in years following drought-reduced yields, and both \textcite{Elmi2004} and \textcite{Elmi2002} note that further reductions could be achieved if controlled drainage is kept operational during early spring (and late fall), when NO$_3$-N losses are most severe and drainage is not needed for crop production.

\subsection{Cropping systems: cover crops, perennials, and rotations}
Cropping-system change is widely promoted but consistently flagged as needing more field evidence, especially across the diversity of soils and climates. \textcite{Kladivko2014} call directly for further research, development, and technical assistance to overcome practical obstacles to cover-crop adoption in the Midwest, for additional field experiments measuring cover-crop effects on NO$_3$ losses in drainage water on multiple soil types across the region, and for using the new data to calibrate and validate RZWQM for additional soils and improve regional estimates of cover-crop NO$_3$ reductions. \textcite{Oliveira2025} identify multiple open questions: how variations in cover-crop biomass, termination timing, and rooting depth interact with drainage dynamics to further affect NO$_3^-$ and dissolved N$_2$O losses, and how variable drainage intensity, fluctuating water tables, and differences in cover-crop biomass and species composition influence dissolved CH$_4$. \textcite{Lagzdins2015} list appropriate mitigation measures needing implementation---annual winter crops, winter cover crops where spring crops are the main crop, and no autumn tillage---to reduce nitrate-N losses to drainage water.

Novel and perennial-based options are also proposed for study. \textcite{Nance2007} highlight the potential of intercropping triticale with red clover for minimizing soil NO$_3$-N and urge that research and policy encourage development of winter triticale to reduce soil NO$_3$-N loss from U.S.\ cropping systems. \textcite{Emmett2022} note that improvements in productivity and land-use efficiency in camelina--soybean relay- and double-cropping systems will be essential to the viability of that sustainable-intensification strategy, and that efforts to improve fall establishment and reduce N surpluses may enable desired reductions in nitrate leaching. \textcite{Randall2001} advocate developing alternative cropping systems containing perennial crops to reduce nitrate losses, and both \textcite{Gowda2007} and \textcite{Nangia2008} suggest that attaining nitrate TMDLs or larger reductions may require alternatives beyond fertilizer management---such as partial conversion of cropland to pasture, changes in tile-drain spacing and/or depth, fall cover crops, or wetland restoration. \textcite{Kaur2024c} add that under variable climate, crop rotations may need modification, or cover crops planted in the winter fallow season to minimize N losses and capture unused soil-profile N. \textcite{Ma2007b} recommend explicitly evaluating controlled drainage, crop rotation, cover crops, and N management for reducing NO$_3$-N loss in drain flow, and \textcite{Malone2019} propose field studies (and paired model testing) of N loss to subsurface drainage with and without corn-stover harvest.

\subsection{Drainage water management and water-table control}
Controlled drainage, shallow drainage, and other water-table management approaches are seen as promising mitigation but require optimization, verification under varied climates, and clearer partitioning of their mechanisms. \textcite{Fang2012} stress that careful management of the timing and depth of gate operation is very important to optimize reduced N loss using controlled drainage. \textcite{Wesström2001} argue that further investigation of controlled drainage/water-table (CWT) system management is needed to optimize both environmental and agricultural benefits, and that water balances for conventional and controlled plots must be established before the reasons behind differing outflow patterns can be separated and quantified. \textcite{Smith2011b} note that continued investigation in years with normal precipitation may clarify whether controlled drainage offers advantages over conventional systems for water conservation, and that further work using shallow drainage under various climatic conditions may reveal its NO$_3$-N reduction potential. \textcite{Madramootoo2024} suggest that controlled-drainage adoption may become more relevant under wetter, warmer climates, and highlight the need for practical water-table control-structure management during wetter springs to mitigate NO$_3$-N concentrations while allowing timely field operations. \textcite{Bakhsh2007} additionally suggest variable-rate controlled drainage may be helpful. \textcite{Schott2017b} identify economic and agronomic questions: whether drainage treatment affects yields (via SEW$_{30}$) during springs with frequent, heavy precipitation, how such conditions affect planting dates, and how the economic risk of reducing NO$_3$ loss with drainage water management can be better evaluated. \textcite{Miller2022} argue that more work is needed to understand the collective impacts of automated drainage water management (ADWM) on the coupled timing of flows and composition of water leaving fields, and how slowing tile flows affects nutrient export, before its full efficacy as a BMP can be assessed. \textcite{Barkle2021} suggest water-table management and edge-of-field bioreactors as mitigation options depending on site conditions.

\subsection{Edge-of-field and in-field treatment technologies}
A distinct research frontier concerns engineered nitrate-removal structures---bioreactors, saturated riparian buffers, constructed wetlands, and filter-based drainage---where design optimization and long-term field performance remain open. \textcite{McEachran2020} state that a key remaining challenge is to develop better ways to predict nitrate removal so that more effective saturated riparian buffers (SRBs) can be designed and implemented; they outline methodological advances including deriving an unsteady form of the governing equation by solving a one-dimensional diffusion equation for the head, computing average linear velocity, and solving nitrate transport numerically, and note that using an optimal-width approximation would help determine effective designs. \textcite{Jaynes2014} propose that if their SRB results can be replicated across landscapes, the practice could prevent substantial NO$_3$ from entering midwestern streams, and that if sequestration is an important removal pathway, biomass harvesting (perhaps as bioenergy) could maintain long-term viability. \textcite{Abdalaal2023} recommend improving how nitrate removal is represented in buffer designs, incorporating more sophisticated removal and hydrology components, coupling design with economic analysis or user objective functions targeting specific annual load-reduction goals, and first determining maximum allowable buffer length before optimizing width to meet load-removal and financial targets.

For bioreactors and combined installations, \textcite{Husk2017} call for further study of controlled drainage with and without a denitrification bioreactor to quantify each component's contribution, note that treating greater water volumes at lower efficiency could improve overall N mass removal (to be tested in performance-optimization studies), stress verification with more on-farm installations, and warn that the suitability of slowing subsurface flow through the bioreactor to filter suspended solids and colloids must be evaluated given clogging potential. \textcite{Vallée2014} encourage planting vegetation and maintaining organic matrices (straw, dead leaves) inside reactive riparian zones (RRZs) to increase effectiveness, since plants could raise efficiency through absorption and contributions to degradation processes. \textcite{Randall2001} propose constructing wetland restoration areas or denitrifying ponds to treat drainage water before it enters ditches or rivers. \textcite{Mendes2021} outline a broad research agenda for constructed treatment systems: thoroughly investigating the driving factors behind performance to support local implementation and nitrogen-removal compliance, clarifying the relationship between water-flow dynamics and nitrogen treatment, elucidating design-parameter effects, conducting long-term field-scale tests of nitrogen biogeochemical processes, investigating volatilization in surface-flow constructed wetlands under anaerobic basic conditions, and developing operations or interventions to enhance winter efficiency and stabilize year-round performance. For filter-based drainage, \textcite{Tao2022} note that over long operation, nitrification and denitrification in the soil and filter cannot be neglected, that long-term field experiments with different filter materials are needed, and that explaining adsorption, desorption, nitrification, and denitrification for different filter materials may be a new area of future exploration.

\subsection{Process-based modeling: outstanding algorithmic and evaluation gaps}
Because field trials cannot cover all combinations of practice, soil, and climate, process-based models are expected to fill the gap---yet the literature documents extensive, still-unresolved modeling needs.

\subsubsection{Soil carbon--nitrogen cycling, plant uptake, and crop growth}
Many authors identify soil N-process representation as the weakest link. \textcite{Smith2020} recommend improving biogeochemical processes such as nitrification, denitrification, fermentation, and mineralization, and using an improved hydrology framework to target denitrification at more appropriate ranges of soil water and oxygen contents. \textcite{Fang2012} attribute high errors in monthly flow-weighted nitrate-concentration predictions to the complexity of crop N uptake, soil N cycling, and transformations among C and N pools, and call for improvement there. \textcite{Hu2007} argue that improving SWAT's simulation of soil N processes, soybean N fixation, and maize-yield simulation could yield more reliable predictions of nutrient-management outcomes in tile-drained watersheds. \textcite{Malone2019} recommend modifying SWAT's nutrient-cycling algorithms to improve immobilization and other processes relevant to corn-stover removal, comparing SWAT and RZWQM at the same site with stover-harvest data, modifying RZWQM's stover-harvest routine to remove greater portions of the upper stalk, and using the modified RZWQM of Fang et al.\ (2017) for more accurate grain and stover N. \textcite{Ma2007} list needed RZWQM improvements: accounting for subsurface lateral flow below tiles between adjacent plots, simulating N fixation rather than supplying all soybean N, improving plant N uptake, incorporating surface roughness/detention storage in runoff, reducing winter nitrification so more soil NH$_4$-N remains at planting, and applying system models to long-term field studies while evaluating all components rather than one. \textcite{Ma2007b} reiterate the need for better simulation of management effects on yield and N mineralization, improved soybean N fixation and plant N uptake, and enhanced post-tillage mineralization, and recommend extending RZWQM to other tile-drained Midwest sites to evaluate hypoxia-alleviating practices. \textcite{Ma2008} add that crop responses to N-application timing need improvement, that use of RZWQM without soil microbial growth needs evaluation for other soils and climates (especially where water is limited and soil organic C differs), and that rigorous evaluation of each soil C/N process and pool against data is required before microbial growth's role can be understood. \textcite{Ma2001} propose a wide-ranging RZWQM agenda: reexamining MSEA datasets with the newest version, introducing variable cracking or effective macroporosity, adding phosphorus soil dynamics for manure management, including mobile organic-matter movement for long-term simulation, studying plant responses to water and N stresses and parameterizing new crops such as wheat, standardizing calibration and database development, evaluating responses to irrigation/fertilization/manure/pesticide management at multiple locations, further evaluating new modules (temperature, freeze/snow heat transport, overland-flow and sediment routing, gas emissions), linking RZWQM to a GIS with economic decision support, and modularizing the model. \textcite{Marjerison2016} report ongoing revision of winter processes and of crop-available-N estimation as suitable datasets emerge, plus work to improve simulation of organic-C mineralization from SOM, residue, and manures and of gaseous N losses in the PNM model. \textcite{Bechtold2007} note that a more robust approach to quantifying grass biomass and N uptake may be needed to simulate C and N dynamics in drained grassland. At the watershed scale, \textcite{Qi2025} recommend developing a more physically based model to represent NO$_3^-$ leaching through tile drainage, pairing outlet NO$_3^-$ loads with soil-moisture and N$_2$O observations to reduce uncertainty in denitrification and N-gas estimates, using ``soft data'' to simulate within-watershed dynamics rather than only outlet loads, adding crop-yield/biomass observations to improve growth-induced N dynamics, and introducing a passive humus pool to address underestimation of organic-N loss via surface runoff.

\subsubsection{Winter, frozen-soil, and snow processes}
Cold-season representation is repeatedly singled out as inadequate. \textcite{Feyereisen2006b} state that better algorithms are needed for soil freezing/thawing and spring-thaw drainage events, that more detailed mechanistic plant-N-uptake representation is needed, and that more detailed snow--freeze--thaw algorithms are still being sought for agronomic systems models. \textcite{Yang2007} call for improving DRAINMOD-N to estimate NO$_3$-N loss more accurately in freezing and thawing seasons and to include inputs such as tillage and intercropping. \textcite{Saadat2020} propose investigating such processes in drained fields to improve DRAINMOD predictions in seasonally frozen regions. \textcite{Jiang2020} note that additional winter data---snow depth, melting, and frost-affected hydraulic conductivity---may help test and improve the model. \textcite{Li2024b} suggest adapting an impedance factor for the RZ-SHAW model's frozen-soil hydraulic conductivity, including soil ice-content measurements in field experiments, implementing a calibratable snow--rain air-temperature threshold, and devoting more attention to new winter management for planting windows and nutrient loss. \textcite{Ahmed2007b} state that continued development and testing of RZWQM for winter conditions is required, along with further testing of ET and plant-water-uptake components, simulation of a wider range of crops for BMP evaluation in sensitive groundwater regions, and more field work to identify coefficient values for the OMNI nutrient sub-model. \textcite{Ahmed2007} similarly recommend improving RZWQM's snow algorithm, incorporating the McBride (1983) water-retention model and Brooks--Corey parameter adjustments for Ontario silt loam, more field work for nutrient-component coefficients, and further testing of ET and plant-water-uptake components. \textcite{Yang2021} indicate it is probably necessary to improve RZWQM2's accuracy in simulating winter runoff-N in winter-wheat cover-crop systems. \textcite{Roy2020} call for research on swelling and shrinkage during infiltration of water into frozen soil. \textcite{Motarjemi2021} require improved winter-crop parameterization, further attention to soil-surface and macropore infiltration processes, evaluation of the Daisy macropore setup where drain-connected macropores exist (with extensive parameter calibration), studies revealing N-loss trends under current practices, and improved management to lower NO$_3$-N in drainage under conventional rotations.

\subsubsection{Macropores, preferential flow, and drain configuration}
\textcite{Smith2020} recommend improving runoff simulation under surface crusting, clay cracking, preferential flow, snow dynamics, and freeze--thaw. \textcite{Stamm2002} argue that more attention should be given to actual drain configuration to complete the picture of factors determining solute transport in drained soils. \textcite{Chelil2022b} propose upgrading SIDRA-RU with a new nitrate-transport model to assimilate nitrate observations, investigating the effect of sampling interval (weekly/bi-weekly versus daily discharge), examining second-order interactions when approximating total-effect Sobol indices, integrating a crop-rotation module if data allow, using water-level, concentration, and flux data to better constrain solutions, and considering time-variant parameter estimation with hourly discharge. \textcite{Jeantet2021} suggest adding reference sites to improve SIDRA-RU robustness, treating the S$_{inter}$ parameter as a variable adjusted to prior-year meteorology and soil parameters, coupling with the PESTDRAIN module for pesticide leaching, and combining SIDRA-RU with a nitrate-leaching module to assess water pollution on French drained plots.

\subsubsection{Drain-design decision support and validation}
\textcite{Moriasi2013} propose using the Revised SWAT model to determine optimum drain depth and spacing that reduce NO$_3$-N losses while optimizing yield, but stress that validation studies of the Hooghoudt and Kirkham tile-drain equations in Revised SWAT are needed before use as a decision-support tool. \textcite{Ale2012} recommend validating a nitrate-load methodology by comparison with N models such as DRAINMOD-NII or RZWQM and by evaluation in additional watersheds where loads from both tile drains and streamflow have been measured. \textcite{Ng2000} note that NO$_3$-N leaching modeling for at least one year is required and that further calibration with consecutive multi-year data may improve prediction for free-drainage treatments. \textcite{Skaggs1995} caution that additional testing for other soils, crops, and locations is needed before routine model application. \textcite{Malone2007} argue that overcoming regression-based limitations and comprehensively quantifying yield and nitrate loading requires process-based models, that simulation results under different conditions could feed a database for economic and risk analysis, and that developing regression equations is a step toward a simple, accurate, objective method to quantify management and climate effects on nitrate loading and yield for a region.

\subsubsection{Groundwater interactions and integrated transport}
Coupling drainage models to groundwater is a distinct need. \textcite{Addab2022} recommend using SWAT-MODFLOW-RT3D for physically based, spatially distributed flow and reactive-transport modeling once a salinity module is incorporated, simulating groundwater storage, flow, and groundwater--tile-drain interactions in a physically based manner, and exploring the balance between improved crop yield in tiled areas and decreased downstream water quality. \textcite{Hoang2017} state that improving SWAT's groundwater module to simulate water-table fluctuation, soil saturation, and groundwater--surface-water exchange is necessary to extend SWAT\_LS validity and realistically represent nitrate retention and removal in riparian zones. \textcite{Hojberg2017} argue that optimal regulation requires considering both spatial variation in N leaching and variation in N reduction across groundwater and surface-water systems. \textcite{Elsayed2023} propose developing integrated mathematical models for nutrient transport in surface-water--groundwater interactions, evaluating climate-change effects, assessing crop-management influences, applying machine-learning algorithms to nutrient-transport datasets, and better understanding governing factors of in-stream sediment and nutrient transport. \textcite{Barkle2021} note that site-specific investigations establishing defensible water balances, or detailed groundwater-flow modeling, are needed to realistically estimate the split between artificial drainage and groundwater flows. \textcite{Yang2008} call for further work to improve NH$_4$-N transport and transformation modeling.

\subsection{Greenhouse-gas tradeoffs and indirect N$_2$O emissions}
A prominent frontier is the coupling between nitrate losses and greenhouse-gas emissions, especially indirect N$_2$O. \textcite{OBrien2022} suggest that continued investigation of the linkages between nitrate leaching and N$_2$O emissions may reveal different dynamics when N application and precipitation are temporally disconnected, and note the need to identify strategies that mitigate losses and adapt to changing precipitation and temperature. \textcite{Sawamoto2003} call for more studies on the NO$_3$--N$_2$O relationship in groundwater, on both direct and indirect N$_2$O emissions under typical conditions, and specifically on indirect N$_2$O emission from subsurface drainage---particularly the NO$_3$-leaching--N$_2$O-leaching relationship important for predicting indirect emission. \textcite{Minamikawa2015} state that field identification is still required and that new or revised calculation methods for indirect N$_2$O associated with leaching and runoff are needed to reflect differences in dominant microbial processes, along with evaluation of the tradeoff between indirect N$_2$O emission and N removal by heterotrophic denitrification across groundwater, surface drainage, river, and estuary. \textcite{Minamikawa2013} note further investigation is required to evaluate groundwater-induced N$_2$O emissions through the soil surface, and \textcite{Minamikawa2010} call for clarifying relations between redox conditions and greenhouse-gas production in the subsoil, for $^{15}$N isotope analysis of nitrate and N$_2$O to locate N$_2$O production, and for understanding production, consumption, and transport of greenhouse gases in profiles with shallow groundwater. \textcite{Smith2019} suggest that future studies of dissolved N$_2$O losses from drainage water may reveal whether they exceed surface emissions, and that further research is needed on the benefits of controlled drainage as a climate-adaptation option in the Atlantic region. \textcite{Oliveira2025} urge exploration of long-term impacts of management on cumulative dissolved greenhouse-gas emissions to inform climate-smart strategies. \textcite{Webb2021} recommend, for artificial-water emissions accounting, examining seasonal impacts on emission factors, developing a process-based model as a likely improvement over revising default EFs (supported by more studies), developing gas-transfer-velocity models across specific artificial-water types, exploring N$_2$O-production thresholds in high-NO$_3$ environments, investigating the DOC/NO$_3$ ratio, building a global database on artificial-water size, density, and distribution (beyond the UK), and targeting field measurements across major agricultural industries. \textcite{Fontes2019} plan to monitor both N$_2$O and NO$_3$-N leaching year-round across all experimental units for better total-N-loss estimates, encourage testing remote-sensing technologies for environmental monitoring, and stress attention to temporal and spatial methodology when evaluating both loss pathways together.

\subsection{Monitoring, tracers, and source apportionment}
Improved measurement and tracing underpin many proposed advances. \textcite{Williams2022} call for higher-frequency measurements (sub-daily to daily), further work on the role of groundwater dynamics in DRP transport (including intra- and inter-event concentration patterns and field differences), more comprehensive methods for apportioning nutrient sources to improve estimates of soil and incidental-fertilizer loads, and further exploration of groundwater dynamics' effect on nutrient delivery to drains. \textcite{Hu2024} recommend quantifying $\delta^{18}$O$_{\mathrm{NO_3,nit}}$ in soils across a broad range of environments to understand its variability under varying thermodynamic and kinetic conditions, and coupling $\delta^{18}$O$_{\mathrm{H_2O}}$ and $\delta^{18}$O$_{\mathrm{NO_3,nit}}$ measurements in evaporating soils to quantify evaporative enrichment. \textcite{Briand2017} argue it is essential to monitor spatial heterogeneity and temporal variability of nitrate concentrations and isotopic compositions in topsoil, drains, and ditches to characterize soil's contribution to river nitrate, with temporal monitoring and multiple tracers needed to explain variations. \textcite{Lagzdins2015} highlight the need for additional research on interactions between drain flow and TN/NO$_3$-N concentration variation, and for more detailed analysis of the interaction between timing and rate of mineral fertilizer and outlet TN/TP concentrations. \textcite{Danalatos2022} ask whether near-term changes in flow-weighted nitrate concentration are predictive of long-term changes in nitrate load and yield. \textcite{Turunen2013} call for modeling and measurement studies to quantify the contribution of groundwater outflow to nutrient loading of surface waters and to study how hydrological changes are reflected in nutrient losses and applied in water-protection design. \textcite{Jokinen2024} recommend that groundwater outflow be considered in future N-loss studies and that fractions other than mineral N be considered when modeling nutrient leaching. \textcite{Barkle2021} add that accurate total-N-export estimates at sites with significant groundwater recharge require measuring all relevant N forms and accounting for the redox status of the underlying groundwater.

\subsection{Scaling, spatial targeting, and thresholds}
Translating site findings to landscapes and defining actionable thresholds constitute another set of open problems. \textcite{Bakhsh2005b} conclude that delineated zones transporting higher NO$_3$-N leaching losses require site-specific management to minimize off-site losses. \textcite{Smith2023} propose finding the NO$_3$-N threshold across a range of P saturation, waterlogging vulnerability, soil textures, organic-carbon availability, and temperature to improve the ability to restrict N and P loss, note that scaling up will require multiple techniques, and suggest combining GIS with hydrological modeling and saturation risk to assess DRP and NO$_3$-N loss from intensifying land. \textcite{Kennedy2012} argue that better understanding the coupled effect of controlling factors and drainage density in agricultural catchments is required for a complete mechanistic synthesis of stream NO$_3^-$ and Cl$^-$ dynamics, and that understanding watershed solute-export responses to increased drainage density requires more quantitative understanding of soil-water and groundwater transport via tile drains. \textcite{Salo2006} suggest identifying critical precipitation periods that create runoff, rather than using annual values, and using the resulting equations to estimate N-leaching risks in, for example, life-cycle analysis.

\subsection{Standards, policy, economics, and cross-cutting tradeoffs}
Finally, several authors point to the need for standards, economic analysis, and explicit tradeoff evaluation. \textcite{Elmi2004} note that developing drain-effluent quality guidelines for NO$_3$-N concentration offers a great challenge, and \textcite{Lagzdins2015} argue that scientifically based standards for water-quality classification in agricultural areas would help define nutrient criteria for agricultural runoff separately from rivers. \textcite{Nash2015d} recommend evaluating whether split N application with managed drainage can further reduce annual NO$_3^-$-N loss through tile drainage from claypan soils, noting it may be possible to further reduce losses from tile-drained claypan soils in the midwestern U.S. Overarching tradeoffs remain unresolved: \textcite{Derby2024} state that the tradeoff of increased NO$_3$-N losses remains unclear and requires further evaluation, and \textcite{Addab2022} likewise call for exploring the balance between crop-yield gains in tiled areas and downstream water-quality losses. \textcite{Gamble2018} emphasize unresolved tradeoffs among N-loss pathways when switching to fertigation of manure. Broader research-and-verification programs are also invoked: \textcite{BAKER1975} describe an early two-part project in which nitrogen was to be applied at 100--300 kg~ha$^{-1}$ to determine the effect of fertilization rate on nitrogen losses, illustrating the long-standing nature of these questions, while \textcite{Dattamudi2020} conclude that better N-management planning must be implemented to reduce nitrate-N contamination in tile drainage. \textcite{Fontes2019} frame the ultimate goal as providing timely, accurate information to enhance the efficiency and sustainability of food production---capturing the applied orientation that unifies this research frontier.